\subsection{Сообщающиеся сосуды}
%%%%%%%%%%%%%%%%%%%%%%%%%%%%%%%%%%%%%%%%%%%%%%%%%%%%
%%%%%%%%%%%%%%%%%%%%%%%%%%%%%%%%%%%%%%%%%%%%%%%%%%%%
\z[{\begin{wrapfigure}[3]{r}{0.10\textwidth}%
		\vspace{-30pt}%
		\begin{flushright}%
 			\includegraphics[width=0.10\textwidth]{4.2.1}%
		\end{flushright}%
	\end{wrapfigure}%
	}]%
{В сообщающихся сосудах находятся ртуть и вода. Высота столба воды $68\,\text{см}$. Какая разница высот уровней ртути установилась в правой и левой трубках?
}
%
{$\Delta h = h \frac{\rho_{\text{в}}}{\rho_{\text{рт}}} = 5 \,\text{см}$}

%%%%%%%%%%%%%%%%%%%%%%%%%%%%%%%%%%%%%%%%%%%%%%%%%%%%
\z{В два открытых сверху цилиндрических сообщающихся сосуда наливают ртуть. Сечение одного из них в два раза больше другого. Широкий сосуд наливают водой до края. На какую высоту $h$ поднимется при этом уровень ртути в другом сосуде? Первоначально уровень ртути был на расстоянии $L$ от верхнего края широкого сосуда.
}
%
{$h = L \frac{2\rho_{\text{в}}}{3\rho_{\text{рт}} - \rho_{\text{в}}}$}

%%%%%%%%%%%%%%%%%%%%%%%%%%%%%%%%%%%%%%%%%%%%%%%%%%%%
\z[{\begin{wrapfigure}[5]{r}{0.22\textwidth}%
		\vspace{-15pt}%
		\begin{flushright}%
 			\includegraphics[width=0.22\textwidth]{4.2.3}%
		\end{flushright}%
	\end{wrapfigure}%
	}]%
{В сосуды, соединённые трубкой с краном, налита вода. Гидростатическое давление в точках $A$ и $B$ равно $p_A = 4\,\text{кПа}$ и $p_B = 1\,\text{кПа}$ соответственно, площади поперечного сечения левого и правого сосудов составляют $S_A = 3\,\text{дм}^2$ и $S_B = 6\,\text{дм}^2$ соответственно. Какое гидростатическое давление установится в точках $A$ и $B$, если открыть кран?
}
%
{$p = \frac{p_A S_A + p_B S_B}{S_A + S_B} = 2 \,\text{кПа}$}

%%%%%%%%%%%%%%%%%%%%%%%%%%%%%%%%%%%%%%%%%%%%%%%%%%%%
\z[{\begin{wrapfigure}[5]{r}{0.22\textwidth}%
		\vspace{-40pt}%
		\begin{flushright}%
 			\includegraphics[width=0.22\textwidth]{4.2.4}%
		\end{flushright}%
	\end{wrapfigure}%
	}]%
{Определите давление воздуха над поверхностью жидкости в точке $A$ внутри закрытого участка изогнутой трубки, если $\rho = 800\,\text{кг/м}^3$, $h = 20\,\text{см}$, $p_0 = 101\,\text{кПа}$, $g = 10\,\text{м/с}^2$. Жидкости плотностями $\rho$ и $2\rho$ друг с другом не смешиваются.
}
%
{$P_A = P_0 - \rho g h = 99.4 \,\text{кПа}$}

%%%%%%%%%%%%%%%%%%%%%%%%%%%%%%%%%%%%%%%%%%%%%%%%%%%%
\z[{\begin{wrapfigure}[5]{r}{0.22\textwidth}%
		\vspace{-20pt}%
		\begin{flushright}%
 			\includegraphics[width=0.22\textwidth]{4.2.5}%
		\end{flushright}%
	\end{wrapfigure}%
	}]%
{Определите давление воздуха над поверхностью жидкости в точке $A$ внутри закрытого участка изогнутой трубки, если $\rho = 800\,\text{кг/м}^3$, $h = 20\,\text{см}$, $p_0 = 101\,\text{кПа}$, $g = 10\,\text{м/с}^2$. Жидкости плотностями $\rho$ и $2\rho$ друг с другом не смешиваются.
}
%
{$L = 350 \,\text{мм}, \quad p = 1160 \,\text{мм рт. ст.}$}

%%%%%%%%%%%%%%%%%%%%%%%%%%%%%%%%%%%%%%%%%%%%%%%%%%%%
\z[{\begin{wrapfigure}[6]{r}{0.22\textwidth}%
		\vspace{+10pt}%
		\begin{flushright}%
 			\includegraphics[width=0.22\textwidth]{4.2.6}%
		\end{flushright}%
	\end{wrapfigure}%
	}]%
{Сосуды, частично заполненные ртутью, над которой находится воздух, сообщаются трубками. Левый верхний сосуд и верхняя трубка открыты в атмосферу. Ртуть по трубкам не перетекает. Найдите давление воздуха в точке $A$, ответ выразите в $\text{мм рт. ст.}$. Определите высоту $L$ столба ртути в верхней трубке. Высота $h = 5\,\text{см}$. Атмосферное давление $p_0 = 760\,\text{мм рт. ст.}$.
}
%
{$h_0 = \frac{3\rho_0 \Delta h - \rho h}{\rho_0} = 155 \,\text{мм}$}

%%%%%%%%%%%%%%%%%%%%%%%%%%%%%%%%%%%%%%%%%%%%%%%%%%%%
\z[{\begin{wrapfigure}[5]{r}{0.22\textwidth}%
		\vspace{+0pt}%
		\begin{flushright}%
 			
\includegraphics[width=0.22\textwidth]{4.2.7}%
		\end{flushright}%
	\end{wrapfigure}%
	}]%
{В сосуде, показанном на рисунке, находится ртуть. Горизонтальные сечения трубок одинаковы. В левую трубку налили воду, высота столба которой $h = 80\,\text{мм}$, а в правую — масло, образовавшее столб некоторой высоты $h_0$. После этого в средней трубке уровень ртути поднялся на $\Delta h = 5\,\text{мм}$. Найдите высоту $h_0$ столба масла, налитого в правую трубку. Плотность воды $\rho = 1000\,\text{кг/м}^3$, масла — $\rho_0 = 800\,\text{кг/м}^3$ и ртути — $\rho_1 = 13600\,\text{кг/м}^3$.
}
%
{$m = \frac{H S \left( \rho_1^2 - \rho_2^2 \right)}{2\rho_2} = 18 \,\text{г}$}